% ============================================
% Supplementary Note 6: Natural Sample-Derived Communities and Stabilization Caveats
% ============================================
\subsection*{\rev{Supplementary Note 6: Natural sample-derived communities and stabilization caveats}}
\label{sec:suppl:note6}

\rev{The natural sample-derived community experiments test coalescence after laboratory stabilization, not coalescence of unfiltered environmental communities. Each environmental inoculum underwent seven serial growth-dilution cycles in defined laboratory media before coalescence. This enrichment phase may pre-select taxa and metabolic strategies that thrive under the culture conditions \citep{Goldford2018}, potentially reducing effective ecological heterogeneity and making natural sample-derived communities more similar to the synthetic consortia.}

\rev{The available post-stabilization 16S data indicate that this filtering was not complete at the ASV level. Natural sample-derived parental communities retained higher richness than the synthetic parental communities ($13.7 \pm 7.2$ versus $9.8 \pm 4.8$ ASVs above the 0.1\% relative-abundance threshold; Supplementary Figs.~22--24). Same-source stabilized parental replicates were also more similar to one another than different-source parental communities at both ASV and Genus levels, and coalesced communities remained closer to their own parental-community pair than to unrelated same-medium parental communities (Supplementary Fig.~44). These analyses rule out complete post-stabilization taxonomic collapse to a single shared endpoint, but they do not quantify how much convergence occurred during stabilization.}

\rev{Because the original environmental inocula were not sequenced before enrichment, we cannot measure pre-to-post taxonomic convergence during the stabilization phase. We also cannot directly test functional convergence because these experiments used 16S community profiling rather than metagenomic, metabolomic, or trait-resolved measurements. Thus, pre-selection could contribute to the convergent coalescence patterns observed between synthetic and natural sample-derived communities, and the natural-community results should be interpreted as evidence that the nutrient-dependent increase in Dominance occurs in taxonomically richer laboratory-stabilized communities, not as evidence for unrestricted generality across unfiltered natural ecosystems. Future work using culture-independent time series or broader environmental sampling could help disentangle laboratory adaptation from intrinsic ecological dynamics.}
